% 2026 CUMCM 参赛论文 —— B 题 无线电干扰源的快速自动定位与清除
%
% 模板: latexstudio/CUMCMThesis (社区维护, 非官方出品)
%       cumcmthesis.cls + cumcm2026.sty 已随仓库提交, 无需额外安装
% 编译: xelatex 连续两遍, 或在项目根目录执行 .\build.ps1
%
% 注意: 不要在此处加 “% !TeX program = xelatex” 魔法注释。
% LaTeX Workshop 读到它会改用内置默认工具, 绕过 .vscode/settings.json
% 里的编译方案, 结果只编译一遍, 交叉引用永远是 ??。
%
% 当前进度: 只完成问题一。问题二、三、四仅留标题占位。
% 摘要与关键词待四问全部完成后统一撰写。

% withoutpreface 去掉封面与编号专用页, 这是 2026 年电子版提交的要求。
% 打印纸质版时改回 \documentclass{cumcmthesis} 以生成承诺书和编号页。
\documentclass[withoutpreface, bwprint]{cumcmthesis}

\usepackage{cumcm2026}

% 算法伪代码环境: 模板未提供, 在此补上。
% algorithm 提供浮动体, algpseudocode 提供 \State \If \For 等语句。
\usepackage{algorithm}
\usepackage{algpseudocode}
\floatname{algorithm}{算法}

% 模板用 \newtheorem 定义了 proof, 会排出「证明 5」这类与定理号对不上的编号。
% 改为不编号, 并在末尾加 QED 记号。
\renewenvironment{proof}{\par\noindent\textbf{证明}\quad}{\hfill$\square$\par\medskip}
\renewcommand{\algorithmicrequire}{\textbf{输入:}}
\renewcommand{\algorithmicensure}{\textbf{输出:}}

% ---------- 绘图 ----------
% 技术路线: Python 只算不画, 结果写入 results/*.csv, 由 PGFPlots 读取绘制。
\usepackage{pgfplots}
\pgfplotsset{compat=1.18}
\usetikzlibrary{arrows.meta, positioning, shapes.geometric, calc}

% viridis 的反转版：本文的等值面图画的都是「越小越好」的目标函数,
% 用反转色标使「亮 = 优」, 与直觉一致。
\pgfplotsset{
  colormap={viridisrev}{indices of colormap={
    \pgfplotscolormaplastindexof{viridis},...,0 of viridis}},
}

\graphicspath{{../figures/}{figures/}}

% 数据目录, 便于统一改路径
\newcommand{\datadir}{../results}

% 图形外部化: 图多到编译明显变慢时再启用, 编译参数需带 -shell-escape
% \usepgfplotslibrary{external}
% \tikzexternalize[prefix=../figures/cache/]

%===================================
\title{无线电干扰源的快速自动定位与清除}
\tihao{B}
\baominghao{xxxxxxxx}        % 报名号
\schoolname{XX 大学}
\membera{队员一}
\memberb{队员二}
\memberc{队员三}
\supervisor{}
\yearinput{2026}
\monthinput{9}
\dayinput{13}
%===================================

\begin{document}

\maketitle

\begin{abstract}
% TODO: 四问全部完成后统一撰写摘要。
% 需回答: 解决什么问题、用了什么模型、核心结果是什么、结论有什么实际意义。
待四问全部完成后撰写。
\end{abstract}

\keywords{待定\quad 待定\quad 待定}

\tableofcontents

\newpage

%======================================================================
\section{问题重述}

\subsection{问题背景}

无线电干扰源的定位与清除是无线电频谱管理的重要任务，通常分为初步监测、区域性排查与
精准定位三个阶段。前两个阶段借助固定监测站、移动监测车或无人机完成大范围信号采集，
将干扰源圈定在一片区域内；受环境遮挡与低功率干扰源的限制，最后的精准定位历来需要技术
人员携带便携式测向机徒步排查，依据场强变化与方向特征逐步逼近。为在危险环境中替代人工
并提高效率，需研制搭载于机器狗的干扰源自动定位清除系统，其核心是一套能在最短时间内
完成“搜索—定位—清除”闭环的策略与算法。

目标区域为半径 $1800$ 米的圆域，取圆心为坐标原点，正东为 $x$ 轴正向，正北为 $y$ 轴正向。
干扰源个数未知，各占互不相同的频道。测向机在检测点测得干扰源的\textbf{示向度}，
即由该点指向干扰源的方位角，受当地电磁环境影响存在不超过 $1^\circ$ 的误差。

\subsection{问题提出}

本文需要解决以下四个问题：

\begin{enumerate}[label=(\arabic*)]
  \item 已知若干个检测点坐标及某干扰源关于这些检测点的示向度，给出采用交会定位法
        计算多边形定位区域直径的算法，并判断以定位区域直径为直径的圆能否覆盖该定位区域。
  \item 已知某全向干扰源在一个检测点处测得的示向度，给出第二个检测点的选择策略，
        并给出第二个检测点的候选区域。
  \item 目标区域内全为全向干扰源、总数在 $10\sim16$ 之间但具体个数未知，
        制定机器狗自动搜索定位及清除的策略与算法，在确保全部清除的前提下使总时间最短。
  \item 目标区域内既有全向干扰源又有定向干扰源，定向干扰源的个数与定向方向均未知，
        在此条件下给出检测与清除的策略及算法。
\end{enumerate}

%======================================================================
\section{问题分析}

\subsection{问题一的分析}

问题一是后续三问的几何基础，其本质是一个\textbf{集员估计}问题而非参数估计问题。

题目给出的误差信息是一个\textbf{确界}：示向度与真实方位角之差落在 $[-1^\circ,1^\circ]$ 内，
而未给出任何分布假设。更关键的是，附录 2 明确指出\textbf{同一地点的电磁环境干扰是固定的，
重复检测不会改变检测误差}。这意味着在同一检测点反复测量不产生任何新信息，
基于独立重复观测的平均化降噪、最小二乘拟合与卡尔曼滤波等方法在此均不适用。

因此，本文采用\textbf{集员估计}（set-membership estimation）的思路\cite{schweppe}：
该方法专为“误差未知但有界”的情形设计，其输出是一个必定包含真值的集合而非点估计，
恰与题目给出的误差信息形式相匹配。具体地，
一次检测所提供的信息不是一条射线，而是一个以检测点为顶点、
半角为 $1^\circ$ 的\textbf{楔形区域}，干扰源必落于其中；多次检测的信息通过\textbf{取交集}
而非加权平均来融合。由于楔形的张角远小于平角，每个楔形可精确表示为两个半平面之交，
于是定位区域归结为若干半平面的交，必为凸多边形，其直径可在顶点对上取到。

第二小问表面上是一个“是否”的判断，实则要求给出\textbf{充要判据}并讨论其紧性。
本文将说明：该问题与最小覆盖圆半径是否达到其下界等价，答案一般为否，
且保证覆盖所需的半径由 Jung 定理给出。

\subsection{问题二的分析}

问题二在问题一的几何框架上增加了一个本质变化：\textbf{检测点由外生给定变为由我们选择}。
问题一的检测点如同已布设的固定监测站，只需评价其定位结果；而问题二要回答“应当选在哪里”，
于是“选”的代价必须进入模型。

分析由三个层次递进。第一，一次检测之后干扰源的可行集是一条纵长 $1495$ 米、
横宽不足 $53$ 米的细长扇形带，方位已定到 $\pm1^\circ$ 而距离几乎一无所知，
因此第二个检测点的作用是\textbf{确定距离}而非改善方位。第二，交会几何存在闭式解，
可证明定位区域直径在“正横站位”处取极小，且\textbf{不存在对全部可能源距一致最优的第二检测点}
（定理 \ref{thm:p2none}）——这说明“使定位区域最小的第二点”这一提法不适定，必须稳健化；
而不同的稳健化准则会给出相距数百米的不同最优点。第三，题目明确只派出\textbf{一只}机器狗，
第二个检测点只能靠其自行走到，行进消耗的正是题目要求最小化的时间，
故“较好的定位效果”不能脱离行程孤立定义。

据此，本文以\textbf{期望总清除时间}为准则，并利用“机器狗速度 $5$ 米/秒、一次检测 $5$ 秒”
这一题设，把时间与路程统一到\textbf{等效行程}量纲上（一次检测等价于 $25$ 米路程），
从而使各项代价可直接相加比较。答案呈现为一个区域而非一个点，
其数学根源是最优侧移量对代价比的\textbf{立方根}依赖；这一自由度也为问题三中
“一次停靠服务多个干扰源”留出了调度空间。

\subsection{问题三、四的分析}

% TODO: 待问题三、四完成后补写。

%======================================================================
\section{模型假设}

在不改变题目核心约束的前提下，作如下假设。其中假设 (1)--(3) 直接来自题目，
假设 (4)--(6) 是本文为完成定量分析所作的建模约定，并非题目直接给定。

\begin{enumerate}[label=(\arabic*)]
  \item 测向机在检测点测得的示向度与干扰源真实方位角之差的绝对值不超过 $1^\circ$
        （题目附录 2 给定）。
  \item 示向度误差由检测点所处的电磁环境决定，在同一地点重复检测不改变该误差
        （题目附录 2 给定）。
  \item 所有干扰源位于半径 $1800$ 米的目标区域内；检测点位置不受此限制
        （题目附件 2 第 1.1 节给定）。
  \item 假设示向度误差 $\varepsilon_{i,c}$ 由检测点位置 $i$ 与频道 $c$ 共同决定，
        对固定的 $(i,c)$ 保持不变，不同 $(i,c)$ 之间不作相关性假设。题目仅说明
        误差由地点决定，未言明同一地点不同频道之间的关系，故此处取不多不少的约定。
  \item 在生成数值算例时，假设误差在 $[-1^\circ,1^\circ]$ 上服从均匀分布。
        题目仅给出误差的确界而未指定分布，均匀分布是仅知支撑区间时的最大熵选择。
        需要强调的是，\textbf{本文问题一的全部结论只依赖确界 $|\varepsilon|\le 1^\circ$，
        与分布形态无关}，该假设仅用于构造算例，其影响将在第 \ref{sec:p1check} 节说明。
  \item 忽略测向机与干扰源的高度差异，视为同一平面上的二维问题（题目附录 2 给定）。
  \item 各干扰源的有效接收半径 $R_{\mathrm{eff}}\in[1000,1500]$ 米，且互不相同、
        无法通过接口查询（题目附录 2 与附件 2 第 2.1 节给定）。因此“必定能收到信号”
        只能按下界 $1000$ 米计算。
  \item 在问题二计算期望值时，假设干扰源在目标区域内均匀分布、$R_{\mathrm{eff}}$ 在
        $[1000,1500]$ 上均匀分布，由此导出源距的先验密度
        $p(\rho)\propto\rho\cdot\Pr(R_{\mathrm{eff}}\ge\rho)$。题目未给出干扰源位置的分布，
        该假设仅用于取期望；其影响将在第 \ref{sec:p2check} 节以敏感性分析说明。
\end{enumerate}

%======================================================================
\section{符号说明}

\begin{table}[H]
  \centering
  \caption{符号说明}
  \label{tab:symbols}
  \begin{tabular}{clc}
    \toprule
    符号 & 含义 & 单位 \\
    \midrule
    $S_i=(x_i,y_i)$ & 第 $i$ 个检测点的坐标，$i=1,2,\ldots,n$ & 米 \\
    $\theta_{i,c}$  & 检测点 $i$ 对频道 $c$ 测得的示向度，取值于 $[0^\circ,360^\circ)$ & 度 \\
    $\varepsilon_{i,c}$ & 检测点 $i$ 对频道 $c$ 的示向度误差 & 度 \\
    $\delta$        & 示向度误差半角，本文取 $\delta=1^\circ$ & 度 \\
    $R$             & 目标区域半径，$R=1800$ & 米 \\
    $G$             & 干扰源的真实位置 & --- \\
    $W_i$           & 检测点 $i$ 确定的楔形区域 & --- \\
    $L=\bigcap_i W_i$ & 定位区域 & --- \\
    $d$             & 定位区域的直径，$d=\max_{P,Q\in L}|PQ|$ & 米 \\
    $(A,B)$         & 实现直径的顶点对，$|AB|=d$ & --- \\
    $R_{\min}$      & 定位区域的最小覆盖圆半径 & 米 \\
    $\varphi$       & 两个检测点对干扰源的交会角 & 度 \\
    $m$             & 定位区域的顶点个数 & --- \\
    $\ell,\ \ell^{*}$ & 检测点到干扰源的距离上限及其可一次清除的阈值（定理 \ref{thm:clearable}） & 米 \\
    $E$             & 顶点超出“以直径为直径的圆”的最大距离 & 米 \\
    \midrule
    \multicolumn{3}{l}{\emph{以下符号用于问题二（第 \ref{sec:p2} 节）}} \\
    $\rho,\ \rho_{\mathrm{hi}}$ & 干扰源到第一检测点的距离 $|S_1G|$ 及其上界 & 米 \\
    $R_{\mathrm{eff}}$ & 干扰源的有效接收半径，$R_{\mathrm{eff}}\in[1000,1500]$ & 米 \\
    $R_c$           & 清除半径，$R_c=20$ & 米 \\
    $a,\,b,\,s$     & 第二检测点的前进量、侧移量与沿轴错位 $s=\rho-a$ & 米 \\
    $T$             & 自第一次检测至该源被清除的总虚拟时间 & 秒 \\
    $\lambda,\mu,\kappa$ & 策略参数，见算法 \ref{alg:p2} & --- \\
    \bottomrule
  \end{tabular}
\end{table}

问题一只涉及单个干扰源，故在第 \ref{sec:p1} 节中略去频道下标 $c$，简记为 $\theta_i$、
$\varepsilon_i$。问题三、四需同时处理多个频道，届时恢复双下标。
此外，定理 \ref{thm:counter} 中的 $\rho,h,r_1,r_2,s$ 为该构型的局部符号，均在定理陈述中就地定义。

%======================================================================
\section{模型的建立与求解}

\subsection{问题一：交会定位区域的直径与覆盖圆问题}
\label{sec:p1}

\subsubsection{定位区域的数学刻画}

设干扰源位于 $G$，检测点 $S_i$ 处测得的示向度为 $\theta_i$。由假设 (1)，真实方位角与
$\theta_i$ 之差不超过 $\delta$，故 $G$ 必落在以 $S_i$ 为顶点、$\theta_i$ 为中心方向、
半角为 $\delta$ 的\textbf{楔形区域}内：

\begin{equation}
  \label{eq:wedge}
  W_i=\Bigl\{P\in\mathbb{R}^2:\ \bigl|\angle(P-S_i)-\theta_i\bigr|\le\delta\Bigr\},
  \qquad
  L=\bigcap_{i=1}^{n}W_i ,
\end{equation}

其中 $\angle(\cdot)$ 表示向量的方位角，$L$ 称为\textbf{定位区域}。当 $n=2$ 时，
$L$ 即题目附录 2 图 2 中的红色四边形；问题一要求将其推广到任意 $n$。

记 $u_i^{\mathrm{lo}}$、$u_i^{\mathrm{hi}}$ 为方向 $\theta_i\mp\delta$ 的单位向量。
由于楔形张角 $2\delta=2^\circ$ 远小于 $180^\circ$，$W_i$ 恰为两个半平面之交：

\begin{equation}
  \label{eq:halfplane}
  W_i=\bigl\{P:\ u_i^{\mathrm{lo}}\times(P-S_i)\ge 0\bigr\}\cap
      \bigl\{P:\ u_i^{\mathrm{hi}}\times(P-S_i)\le 0\bigr\},
\end{equation}

$\times$ 为二维叉积。于是 $n$ 次检测对应 $2n$ 个半平面，$L$ 为其交集。半平面是闭凸集，
闭凸集之交仍为闭凸集，且每个半平面至多贡献一条边，故\textbf{$L$ 为凸多边形，
边数至多 $2n$}。

$2n$ 只是上界。若某检测点距干扰源较远、或其示向方向与已有约束相近，
相应半平面即成为\textbf{冗余约束}——第 \ref{sec:p1num} 节的算例 C 用了 $5$ 个检测点，
定位区域却只有 $6$ 个顶点，即 $4$ 个半平面冗余。可见\textbf{决定定位精度的是检测点的
几何布置而非个数}，此结论直接指向问题二的选点策略。

\subsubsection{直径的计算与算法}

定位区域的直径定义为 $d=\max_{P,Q\in L}|PQ|$。固定 $P$ 时 $Q\mapsto|PQ|$ 为凸函数，
凸函数在凸多边形上的最大值必在顶点取到；对 $P$ 再作同样论证，即得

\begin{equation}
  \label{eq:diam}
  d=\max_{1\le a<b\le m}\lvert V_aV_b\rvert ,
\end{equation}

$V_1,\ldots,V_m$ 为 $L$ 的顶点。故只需枚举顶点对，无须考虑边上与内部的点。据此得算法 \ref{alg:diam}。

\begin{algorithm}[H]
\caption{定位区域与直径的计算}
\label{alg:diam}
\begin{algorithmic}[1]
\Require 检测点 $S_i$ 与示向度 $\theta_i$（$i=1,\ldots,n$），误差半角 $\delta$
\Ensure 定位区域顶点表 $V$，直径 $d$，直径点对 $(A,B)$
\If{$\bigcap_{i=1}^{n}[\theta_i-\delta,\ \theta_i+\delta]\ne\varnothing$}
  \State \Return “定位区域无界” \Comment{有界性判据，证明见附录 \ref{app:proof}}
\EndIf
\State $P\leftarrow$ 一个足够大的辅助方框 \Comment{数值脚手架，与目标区域无关}
\For{$i=1$ \textbf{to} $n$}
  \State $P\leftarrow\textsc{Clip}\bigl(P,\ \{X: u_i^{\mathrm{lo}}\times(X-S_i)\ge 0\}\bigr)$
  \State $P\leftarrow\textsc{Clip}\bigl(P,\ \{X: u_i^{\mathrm{hi}}\times(X-S_i)\le 0\}\bigr)$
  \If{$P=\varnothing$} \State \Return “读数矛盾，无可行定位区域” \EndIf
\EndFor
\State $V\leftarrow$ 去除 $P$ 中的重合点与共线冗余点
\State 枚举顶点对求 $d$ 与 $(A,B)$，返回 $V,\ d,\ (A,B)$
\end{algorithmic}
\end{algorithm}

\textsc{Clip} 为 Sutherland--Hodgman 半平面裁剪。第 4 行的辅助方框仅为使裁剪从有界多边形
开始，本文取边长 $10R$，\textbf{与目标区域无关}；此处不宜用目标区域的外接正方形，
否则定位区域跨出圆域时会被错误裁剪并切出并不存在的直边。

\paragraph{复杂度与退化情形.}
裁剪为 $O(nm)$，求直径为 $O(m^2)$，$m\le 2n$。实际 $n$ 不超过数十，代价可忽略；
若 $n$ 极大可用旋转卡壳将求直径降至 $O(m)$。平面上凸集之交只有五种形态，
构成完整分类，见表 \ref{tab:degenerate}。单点与线段虽为测度零事件，
但直径均有定义，算法照常计算而不作特例跳过。

\begin{table}[H]
  \centering
  \caption{定位区域的五种形态及其处理}
  \label{tab:degenerate}
  \begin{tabular}{lll}
    \toprule
    形态 & 判定 & 直径与处理 \\
    \midrule
    空 & 裁剪后无顶点 & 无定义；读数矛盾，检出并报错 \\
    无界 & 各示向度区间之交非空（含 $n=1$） & $+\infty$；报无界并提示增大基线 \\
    单点 / 线段 & 约束恰好相切 / 共线 & $0$ / 线段长；照常计算 \\
    多边形 & 一般情形 & 顶点对最大距离 \\
    \bottomrule
  \end{tabular}
\end{table}

无界并非算法失效，而是一个有物理意义的诊断。由于 $\theta_i$ 是 $S_i$ \emph{指向} $G$ 的
方位角，$\theta_1-\theta_2$ 恰等于 $G$ 处的交会角 $\varphi$，故 $n=2$ 时

\begin{equation}
  \label{eq:unbounded2}
  L\ \text{无界}\iff|\theta_1-\theta_2|\le2\delta\iff\varphi\le2\delta=2^\circ .
\end{equation}

即两次检测的方向分辨率不足以区分，第二次检测未提供任何新的方向信息。
此时正确的应对是\textbf{增大基线后重测}，而非继续解算。当两检测点到干扰源距离相近且间距 $b$ 远小于该距离 $r$ 时 $\varphi\approx b/r$，
故千米量级上两检测点相距不足 $35$ 米即会触发；
该情形在问题三、四中机器狗边走边测时并不罕见。条件 $\varphi>2\delta$ 亦构成问题二中
第二检测点候选区域的必要约束。

\subsubsection{覆盖圆问题}
\label{sec:cover}

设直径由顶点对 $(A,B)$ 达到。由泰勒斯定理，顶点 $V$ 位于以 $AB$ 为直径的圆内或圆上
当且仅当 $\angle AVB\ge 90^\circ$，写成内积形式即得判据

\begin{equation}
  \label{eq:criterion}
  \text{该圆覆盖 }L\iff (A-V)\cdot(B-V)\le 0,\quad \forall V\in V(L).
\end{equation}

判据只需一次点积，无须反三角函数，且返回的越界顶点即反例点。\footnote{题目将“定位区域直径”
定义为一个数值，故“以定位区域直径为直径的圆”亦可理解为“直径等于该数值、位置任意的圆”。
两种理解等价：因任何覆盖 $L$ 的半径 $r$ 的圆满足 $d\le 2r$，故 $R_{\min}\ge d/2$ 恒成立；
若存在直径为 $d$ 的覆盖圆则 $R_{\min}=d/2$，此时 $A,B$ 距离恰等于该圆直径，
只能是某条直径的两端，圆心必为 $AB$ 中点。}

\paragraph{答案：一般不能.}
对一般平面凸集，等边三角形即为反例（直径 $1$，最小覆盖圆半径 $1/\sqrt3>0.5$）。
但定位区域是楔形之交这一特殊类型，故需针对本问单独证明。下述定理给出一个可完全解析求解的构型。

\begin{theorem}[三点均布构型的解析反例]
\label{thm:counter}
设三个检测点位于以干扰源 $G$ 为圆心、半径 $\rho>0$ 的圆周上且互成 $120^\circ$，
示向度均无误差，楔形半角 $\delta\in(0^\circ,15^\circ]$。记
\begin{equation}
  \label{eq:r1r2}
  h=\rho\sin\delta,\quad
  r_1=\frac{h}{\cos(30^\circ-\delta)},\quad
  r_2=\frac{h}{\cos(30^\circ+\delta)},\quad
  s=\frac{r_2-r_1}{2}>0 .
\end{equation}
则有：
\begin{enumerate}[label=(\alph*)]
  \item $L$ 为等边六边形，顶点方向相隔 $60^\circ$、到 $G$ 的距离在 $r_1$ 与 $r_2$ 间交替，
        故 $L$ 只具 $3$ 次旋转对称性，\textbf{并非正六边形}；
  \item 直径 $d=r_1+r_2$ 由一对对径顶点取到，且以 $AB$ 为直径的圆\textbf{不能}覆盖 $L$，
        最大超出量 $E=\sqrt{r_2^2+r_2s+s^2}-(r_2-s)>0$；
  \item 该构型的覆盖比为 $R_{\min}/(d/2)=1+\tan\delta/\sqrt3$，
        在 $\delta=15^\circ$ 处恰等于 $2/\sqrt3$，即\textbf{精确达到 Jung 界}。
\end{enumerate}
\end{theorem}

\begin{proof}
\textbf{(1) 六条边界线同切于一圆。} $W_k$ 的两条边界射线过 $S_k$、与 $S_kG$ 方向成
$\pm\delta$ 角，故 $G$ 到每条边界线的距离均为 $|S_kG|\sin\delta=\rho\sin\delta=h$。

\textbf{(2) 法向夹角交替。} 设 $S_k$ 方位角为 $\alpha_k=\alpha_0+120^\circ k$，
直接计算得 $W_k$ 两条边界线的外法向角为 $\alpha_k\mp(90^\circ-\delta)$。
六个法向角排序后，相邻之差\textbf{交替}为 $60^\circ-2\delta$ 与 $60^\circ+2\delta$，
故 $L$ 只具 $3$ 次旋转对称性——这正是 $L$ 非正六边形的根源。

\textbf{(3) 顶点半径。} 两条距 $G$ 为 $h$、外法向夹角为 $\Delta$ 的切线，其交点位于两法向的
角平分线上；设该点距 $G$ 为 $t$，则法向投影给出 $t\cos(\Delta/2)=h$，即 $t=h/\cos(\Delta/2)$。
代入 $\Delta=60^\circ\mp2\delta$ 即得式 \eqref{eq:r1r2} 的 $r_1$、$r_2$，且顶点方向相隔
$60^\circ$、半径交替。由 $0<\delta<30^\circ$ 得 $\cos(30^\circ+\delta)<\cos(30^\circ-\delta)$，
故 $r_2>r_1$，$s>0$，(a) 得证。

\textbf{(4) 直径由对径顶点对取到。} 顶点方向均布且半径交替，故对径顶点半径必不相同，
距离为 $r_1+r_2$；与其余顶点对比较可知，$\delta\in(0^\circ,15^\circ]$ 时该距离最大
（比较过程见附录 \ref{app:proof}），故 $d=r_1+r_2$。

\textbf{(5) 覆盖判定。} 设 $A$（半径 $r_2$，方向 $\psi$）与 $B$（半径 $r_1$，方向
$\psi+180^\circ$）为直径点对，则圆心 $C=s\hat\psi$、半径 $R=(r_1+r_2)/2=r_2-s$。
取方向 $\psi+120^\circ$、半径 $r_2$ 的顶点 $V$，由余弦定理
$|VC|^2=r_2^2+s^2-2r_2s\cos120^\circ=r_2^2+r_2s+s^2$，故

\begin{equation}
  \label{eq:key}
  |VC|^2-R^2=\bigl(r_2^2+r_2s+s^2\bigr)-\bigl(r_2-s\bigr)^2=3r_2s
  =\frac{3\rho^2\sin^3\delta}{2\cos^2(30^\circ+\delta)\cos(30^\circ-\delta)}>0 ,
\end{equation}

即 $V$ 严格落在圆外，(b) 得证。(c) 的推导见附录 \ref{app:proof}。
\end{proof}

式 \eqref{eq:key} 是\textbf{精确代数恒等式}，其符号不依赖任何数值计算。
该定理另有两点值得注意：其一，结论\textbf{与 $\rho$ 无关}，即检测点距离多远都同样失效，
并非刻意构造的特例；其二，由 \eqref{eq:key} 知 $E$ 随 $s$ 单调，而 $\delta\to0$ 时
$r_1\to r_2$、$s\to0$、$E\to0$——此时楔形退化为等宽条带、六边形成为正六边形、
直径圆恰为其外接圆，故\textbf{正六边形恰是“能覆盖”的临界情形，是楔形的锥度把它推过了
临界线}。图 \ref{fig:excess} 给出 $E$ 随 $\delta$ 的变化。

\begin{figure}[H]
  \centering
  \begin{tikzpicture}
    \begin{axis}[
      width=0.62\linewidth, height=4.6cm,
      xlabel={楔形半角 $\delta$ / 度}, ylabel={超出量 $E$ / 米},
      xmin=0, xmax=2.2, ymin=0, ymax=1.2,
      grid=major, tick label style={font=\scriptsize}, label style={font=\small},
      legend style={font=\scriptsize, at={(0.03,0.97)}, anchor=north west, draw=black!30},
    ]
      % 由 E = sqrt(r2^2+r2*s+s^2)-(r2-s) 解析式计算, rho=900 米
      \addplot[thick, blue!70!black, domain=0.001:2.2, samples=120] {
        sqrt( (900*sin(x)/cos(30+x))^2
            + (900*sin(x)/cos(30+x))*((900*sin(x)/cos(30+x)-900*sin(x)/cos(30-x))/2)
            + ((900*sin(x)/cos(30+x)-900*sin(x)/cos(30-x))/2)^2 )
        - ( 900*sin(x)/cos(30+x) - (900*sin(x)/cos(30+x)-900*sin(x)/cos(30-x))/2 ) };
      \addlegendentry{$\rho=900$ 米}
      \addplot[only marks, mark=*, mark size=1.8pt, red] coordinates {(1,0.2749)};
      \addlegendentry{本题 $\delta=1^\circ$}
    \end{axis}
  \end{tikzpicture}
  \caption{超出量 $E$ 随楔形半角 $\delta$ 的变化。$\delta\to0$ 时 $E\to0$，
           即正六边形是覆盖与否的临界情形；本题 $\delta=1^\circ$ 处 $E=0.2749$ 米}
  \label{fig:excess}
\end{figure}

$n=2$ 亦可失效，机理相同：楔形两边界射线以 $2\delta$ 发散，两检测点等距且误差对称时
锥度抵消、四边形退化为平行四边形而恰好被覆盖，一旦对称性被破坏即可能越界。
第 \ref{sec:p1num} 节的算例 A 即为一例。

\paragraph{保证覆盖所需的半径与清除判据.}
由 Jung 定理\cite{jung}，平面上直径为 $d$ 的点集必可被半径 $d/\sqrt3$ 的圆覆盖，
等边三角形取等号，故

\begin{equation}
  \label{eq:jung}
  \frac{d}{2}\ \le\ R_{\min}\ \le\ \frac{d}{\sqrt3}\approx 0.577\,d .
\end{equation}

题目规定光学探测仪须在距干扰源 $20$ 米内方可清除。若要求\textbf{站在单一位置一次完成清除}
且保证成功，则定位区域须整体落入半径 $20$ 米的圆内，故准确判据为 $R_{\min}\le 20$ 米
（站位取最小覆盖圆圆心）；由式 \eqref{eq:jung}，$d\le 20\sqrt3\approx34.6$ 米是便于使用的
\textbf{充分条件}，而 $d\le 40$ \textbf{并不充分}。另一方面，$\pm1^\circ$ 误差在距离 $r$ 处的
横向不确定度全宽为 $2r\tan1^\circ$，$r=1000$ 米时即达 $34.91$ 米，与上述阈值相当。
这说明\textbf{千米量级上一次交会的精度恰处于“够与不够”的临界}，
问题三、四必须采用“先粗定位、逼近后补测”的策略。

那么“逼近”到什么程度才够？下述定理给出定量答案，它同时也是问题三、四中
“确保所有干扰源被清除”这一要求的可行性依据。

\begin{theorem}[可清除性保证]
\label{thm:clearable}
设机器狗已接收到干扰源 $G$ 的信号。若其在距 $G$ 均不超过 $\ell$ 的两处取检测点、
且二者相对 $G$ 的张角为 $90^\circ$，则所得定位区域满足
\begin{equation}
  \label{eq:clearbound}
  R_{\min}\ \le\ \frac{\sqrt2\,\ell\sin\delta}{1-2\sqrt2\sin\delta}.
\end{equation}
于是只要
\begin{equation}
  \label{eq:ellstar}
  \ell\ \le\ \ell^{*}=\frac{20\,(1-2\sqrt2\sin\delta)}{\sqrt2\sin\delta}\approx 770\ \text{米}
  \qquad(\delta=1^\circ),
\end{equation}
即有 $R_{\min}\le20$ 米，机器狗在最小覆盖圆圆心处一次清除必定成功。
\end{theorem}

\begin{proof}
记 $d$ 为定位区域 $L$ 的直径，两检测点距 $G$ 分别为 $r_1,r_2\le\ell$。
对任意 $P\in L\subseteq W_i$，$P-S_i$ 与中心射线方向的夹角不超过 $\delta$，
故 $P$ 到该射线的垂距不超过 $|P-S_i|\sin\delta$；又 $G\in L$，
故 $|P-S_i|\le|P-G|+|G-S_i|\le d+r_i$。于是 $L$ 含于以中心射线为轴、
半宽 $h_i=(d+r_i)\sin\delta$ 的条带内。两条中心射线的夹角等于张角 $90^\circ$，
故 $L$ 含于半宽为 $h_1,h_2$ 的矩形之内，该矩形的外接圆半径与直径分别为
$\sqrt{h_1^2+h_2^2}$ 与 $2\sqrt{h_1^2+h_2^2}$。由 $h_i\le(d+\ell)\sin\delta$ 得
\[
  d\le 2\sqrt2\,(d+\ell)\sin\delta
  \ \Longrightarrow\
  d\le\frac{2\sqrt2\,\ell\sin\delta}{1-2\sqrt2\sin\delta},
\]
代回 $R_{\min}\le\sqrt{h_1^2+h_2^2}\le\sqrt2\,(d+\ell)\sin\delta$ 即得式 \eqref{eq:clearbound}。
\end{proof}

定理 \ref{thm:clearable} 的意义在于 $\ell^{*}\approx770$ 米\textbf{小于有效接收半径的下界
$1000$ 米}：机器狗总能移动到距干扰源 $770$ 米以内，且在该处必定仍能接收到信号，
因而上述两次检测总可执行。当 $G$ 接近目标区域边界时，所需检测位置可能落在
半径 $1800$ 米的圆域之外——题目附件 2 第 1.1 节明确允许机器狗提交目标区域外的位置，
故结论对边界附近的干扰源同样成立。此外，式 \eqref{eq:clearbound} 中的界是保守的：
$\ell=770$ 米时实测 $R_{\min}=19.0$ 米，约为上界的 $95\%$；
且张角在 $60^\circ\sim120^\circ$ 范围内、$\ell=500$ 米时 $R_{\min}$ 均不超过 $17.5$ 米，
故实际执行时无须精确取到 $90^\circ$。

\subsubsection{数值算例与结果}
\label{sec:p1num}

题目未提供问题一的数据，本文自行构造三组算例，生成脚本固定随机种子 $20260910$，
可完全复现。各检测点距干扰源均超过 $5$ 米（否则按附录 2 第 9 条信号过强而无示向度），
且互不重合（否则按假设 (2) 须共用同一误差）。求解结果见表 \ref{tab:p1result}。

\begin{table}[H]
  \centering
  \caption{三组算例的设置与求解结果}
  \label{tab:p1result}
  \begin{tabular}{clcccccc}
    \toprule
    算例 & 布置 & $n$ & $m$ & $d$/米 & $R_{\min}$/米 & 圆能否覆盖 & 超出量 $E$/米 \\
    \midrule
    A & 交会角 $90^\circ$，两点各距源 $1000$ 米 & 2 & 4 & 49.3796 & 24.6900 & 否 & 0.0874 \\
    B & 三点均布 $120^\circ$，误差取零 & 3 & 6 & 36.2834 & 18.3245 & 否 & 0.2749 \\
    C & 五点随机布置 & 5 & 6 & 27.7060 & 13.8530 & 是 & 0 \\
    \bottomrule
  \end{tabular}
\end{table}

算例 A 的检测点布置是预先设定的，两个误差 $-0.9864^\circ$ 与 $+0.8110^\circ$ 由固定
随机种子下的均匀分布生成、\textbf{未针对覆盖判定作任何挑选}，其定位区域仍有一个顶点
超出圆 $0.0874$ 米，判据侧亦有 $(A-V)\cdot(B-V)=+4.324>0$。算例 B 是定理 \ref{thm:counter}
在 $\rho=900$ 米、$\delta=1^\circ$ 处的实例，其定位区域如图 \ref{fig:caseB}；
六条边等长（均 $18.1444$ 米）但外接半径在两值间交替，仅具 $3$ 次旋转对称性，
为\textbf{等边而非正}的六边形，正是定理 \ref{thm:counter} 的结论。算例 C 表明
覆盖失败并非必然。三组算例的定位区域均包含干扰源真实位置，与假设 (1) 一致。

\begin{figure}[H]
  \centering
  \begin{tikzpicture}
    \begin{axis}[
      width=0.62\linewidth, height=6cm,
      axis equal image,
      xlabel={$x$ / 米}, ylabel={$y$ / 米},
      xmin=-26, xmax=26, ymin=-26, ymax=26,
      grid=major,
      legend style={font=\scriptsize, at={(1.03,0.5)}, anchor=west, draw=black!30},
      tick label style={font=\scriptsize}, label style={font=\small},
    ]
      \addplot[thick, draw=blue!70!black, fill=blue!12]
        table [x=x, y=y, col sep=comma] {\datadir/问题一_算例B_定位区域.csv};
      \addlegendentry{定位区域}
      % 圆心与半径取自 results/问题一_求解汇总.csv
      \addplot[red, thick, domain=0:360, samples=181, smooth]
        ({0.1583+18.1417*cos(x)}, {-0.0914+18.1417*sin(x)});
      \addlegendentry{以直径为直径的圆}
      \addplot[green!50!black, thick, dotted, domain=0:360, samples=181, smooth]
        ({18.3245*cos(x)}, {18.3245*sin(x)});
      \addlegendentry{最小覆盖圆}
      \addplot[red, thick, mark=*, mark size=1.4pt]
        coordinates {(15.869,-9.162) (-15.553,8.979)};
      \addlegendentry{直径 $AB$}
      \addplot[only marks, mark=*, mark size=2.6pt, orange] coordinates {(0,18.3245)};
      \addlegendentry{落在圆外的顶点}
    \end{axis}
  \end{tikzpicture}
  \caption{算例 B 的定位区域与两个圆。两圆半径相差仅 $0.18$ 米、图中几乎重合，
           顶部顶点超出以直径为直径的圆 $0.2749$ 米，具体数值见表 \ref{tab:p1result}}
  \label{fig:caseB}
\end{figure}

\subsubsection{模型检验}
\label{sec:p1check}

问题一不存在可供比对的标准答案，故以\textbf{题目给出的定义本身作为判据}，
从三个层次检验算法实现，各层均\textbf{不调用被检验的代码}，以免同一处错误被两侧同时继承。
其一，在检测点周边随机取 $2\times10^5$ 点，比对“依定义由方位角判定”与“依式
\eqref{eq:halfplane} 的不等式判定”；其二，按内法向形式独立推导半平面，
交由 Qhull 算法\cite{qhull}求交，与本文所得顶点逐一配对；其三，在以定位区域为中心、
半宽为其直径三倍的方框内均匀采样 $4\times10^5$ 点，比对“依式 \eqref{eq:wedge} 判定”与
“依所求多边形判定”。三组算例中，第一、三层均无分歧，第二层顶点最大偏差不超过
$2.3\times10^{-12}$ 米。此外采用\textbf{变异测试}向算法注入半角取值错误、
楔形顶点偏移等典型缺陷，各层检验均能检出，表明验证方案确具检错能力。

定理 \ref{thm:counter} 为算法提供了独立于数值计算的校验基准。将 $\rho=900$ 米、
$\delta=1^\circ$ 代入式 \eqref{eq:r1r2}，解析值与算法数值解在有效数字范围内完全一致：

\begin{table}[H]
  \centering
  \caption{算例 B 的解析值与数值解对照}
  \label{tab:caseBcheck}
  \begin{tabular}{lcccc}
    \toprule
     & $r_1$/米 & $r_2$/米 & $d$/米 & 超出量 $E$/米 \\
    \midrule
    解析值   & 17.958851905 & 18.324504190 & 36.283356095 & 0.274919837 \\
    数值解   & 17.958851905 & 18.324504190 & 36.283356095 & 0.274919837 \\
    \bottomrule
  \end{tabular}
\end{table}

该算例坐标量级约 $18$ 米，双精度浮点绝对误差约 $4\times10^{-15}$ 米，
与超出量 $0.2749$ 米相差 $13$ 个数量级，可排除数值误差导致覆盖判定失误的可能。

\paragraph{对假设的灵敏度.}
均匀分布假设 (5) 仅用于生成算例。本节全部结论——包括定位区域的定义与结构、直径与覆盖判据、
定理 \ref{thm:counter} 与 \ref{thm:clearable}、表 \ref{tab:degenerate} 的形态分类——
\textbf{只依赖确界 $|\varepsilon_i|\le\delta$，与误差的分布形态及相关结构无关}。这是集员估计相较于
最小二乘、卡尔曼滤波等方法的固有优势：题目仅给出误差上界时，前者无须额外分布假设
即可给出严格的包含性结论。

作为补充，本文以干扰源置于原点、$n$ 个检测点距其 $300\sim1200$ 米、方位与误差均匀随机的
方式各扫描 $4000$ 组配置，考察比值 $R_{\min}/(d/2)$ 的上确界：$n=2$ 时为 $1.0005$
（覆盖失败率 $1.0\%$），$n=3$ 时为 $1.1308$（$17.2\%$），$n=4$ 时为 $1.1383$（$26.4\%$），
后者已达 Jung 界的 $98.6\%$。该结果为数值扫描所得，不构成证明，
但与定理 \ref{thm:counter}(c) 中“Jung 界可被楔形之交精确取到”的解析结论相互印证。

\subsection{问题二：第二个检测点的选择策略}
\label{sec:p2}

\subsubsection{一次检测后的可行集与评价准则}
\label{sec:p2set}

设第一个检测点为 $S_1$、测得示向度为 $\theta_1$，记 $\rho=|S_1G|$。除楔形 $W_1$ 外，
题目还给出两条信息：既然在 $S_1$ 收到了信号，则 $\rho\le R_{\mathrm{eff}}\le1500$ 米；
既然返回的是示向度而非“距离过近”，则 $\rho>5$ 米。再计入干扰源必在目标区域内，得可行集
\begin{equation}
  \Omega_1 = W_1 \cap \{\,5<\rho\le\rho_{\mathrm{hi}}\,\}\cap D(O,R),
  \qquad
  \rho_{\mathrm{hi}}=\min\bigl(1500,\ \text{沿楔形至目标圆边界的距离}\bigr).
  \label{eq:p2omega}
\end{equation}

$S_1$ 位于原点时 $\rho_{\mathrm{hi}}=1500$ 米，$\Omega_1$ 纵长 $1495$ 米而最大横宽仅
$2\rho_{\mathrm{hi}}\tan\delta=52.4$ 米，形如一条细长的扇形带（图 \ref{fig:p2set}(a)）。
可见\textbf{方位已定到 $\pm1^\circ$，而距离几乎一无所知}，于是：

\begin{center}
\emph{第二个检测点的作用不是改善方位精度，而是确定距离。}
\end{center}

“较好的定位效果”需要一个可计算的判据。定位区域的\textbf{面积}是一个自然的候选，
但对本题并不适用：由附录 2 第 8 条，清除要求站位与干扰源相距不超过 $R_c=20$ 米，
故可否一次清除取决于定位区域在\textbf{最坏方向}上的尺度，而面积对形状不敏感。
表 \ref{tab:p2area} 给出一组面积几乎相同、可清除性却相反的实例。

\begin{table}[H]
  \centering
  \caption{面积相同而可清除性相反的两个第二检测点（$\rho=1000$ 米）}
  \label{tab:p2area}
  \begin{tabular}{lccccc}
    \toprule
    第二检测点 $(a,b)$ & 面积/米$^2$ & 直径 $d$/米 & $R_{\min}$/米 & 交会角 $\varphi$ & 可否一次清除 \\
    \midrule
    $(800,\,200)$ & 488.6 & 60.16 & 30.08 & $45.0^\circ$ & 否 \\
    $(950,\,395)$ & 489.2 & 39.49 & 19.75 & $82.8^\circ$ & 是 \\
    \bottomrule
  \end{tabular}
\end{table}

这并非巧合。以 $S_1$ 为原点、$\theta_1$ 为 $x$ 轴正向建立局部标架，
记第二检测点为 $S_2=(a,b)$（$a$ 为沿示向线的前进量，$b$ 为侧移量），
干扰源为 $G=(\rho,0)$，并记沿轴错位 $s=\rho-a$，参见图 \ref{fig:p2set}(b)。
则两楔形之交的面积为 $4\rho\tan^2\!\delta\,(s^2+b^2)/b$，对 $b$ 取极小得 $b=|s|$，
即交会角 $\varphi=45^\circ$——恰为表 \ref{tab:p2area} 中不可清除的一行。
\textbf{面积准则会系统性地选中错误的构型}。
故本文沿用问题一的口径，以直径 $d$ 度量定位效果，以 $R_{\min}\le R_c$ 作为可一次清除的充要判据。
对 $n=2$，二者的关系比一般情形更紧密：

\begin{lemma}[两楔形之交的覆盖性]
\label{lem:p2para}
若两楔形之交未被式 \eqref{eq:p2omega} 中的距离约束截断，则该交集为平行四边形，
且 $R_{\min}=d/2$ 严格成立。
\end{lemma}

\begin{proof}
在干扰源邻域内，楔形 $W_i$ 可视为宽度 $2r_i\tan\delta$ 的条带，两条带之交为平行四边形。
其较长对角线所对的内角为钝角，由泰勒斯定理该顶点落在以长对角线为直径的圆内，另一顶点同理；
故该圆覆盖全部顶点。又恒有 $R_{\min}\ge d/2$，两者相等。
\end{proof}

$2\times10^4$ 组随机构型中仅 $2.94\%$ 因被距离上界截断而不满足，且此时 $R_{\min}/(d/2)$
最大仅 $1.000577$。因此在问题二范围内，判据“$d\le 2R_c=40$ 米”与“可一次清除”基本等价，
这与问题一 $n\ge3$ 时的情形形成对照。

\begin{figure}[H]
  \centering
  \begin{tikzpicture}
    \begin{axis}[
      name=left, width=0.46\linewidth, height=4.2cm,
      xlabel={沿示向线方向 / 米}, ylabel={横向 / 米},
      xmin=-60, xmax=1580, ymin=-58, ymax=58,
      grid=major, title={(a) 可行集 $\Omega_1$（横向放大 $14$ 倍）},
      title style={font=\small}, tick label style={font=\scriptsize},
      label style={font=\scriptsize},
    ]
      \addplot[draw=blue!70!black, fill=blue!15, thick]
        table [x=x, y=y, col sep=comma] {\datadir/问题二_可行集.csv} \closedcycle;
      \addplot[only marks, mark=triangle*, mark size=2.4pt, black] coordinates {(0,0)};
      \node[font=\scriptsize] at (axis cs:120,-38) {$S_1$};
      \node[font=\scriptsize] at (axis cs:760,34) {纵长 $1495$ 米，最大横宽 $52$ 米};
    \end{axis}
    \begin{axis}[
      name=right, at={($(left.east)+(1.4cm,0)$)}, anchor=west,
      width=0.40\linewidth, height=4.2cm, axis equal image,
      xlabel={$x$ / 米}, ylabel={$y$ / 米},
      xmin=-80, xmax=1150, ymin=-120, ymax=330,
      grid=major, title={(b) 第二检测点的参数},
      title style={font=\small}, tick label style={font=\scriptsize},
      label style={font=\scriptsize},
    ]
      % 两条示向线均须过 G=(900,0)：S1 的沿 x 轴，S2 的方向为 (600,-150)
      \addplot[gray, thick] coordinates {(0,0) (1100,0)};
      \addplot[brown!70!black, thick] coordinates {(300,150) (1100,-50)};
      \addplot[only marks, mark=triangle*, mark size=2.4pt, black] coordinates {(0,0)};
      \addplot[only marks, mark=square*, mark size=2.0pt, black] coordinates {(300,150)};
      \addplot[only marks, mark=star, mark size=3.0pt, red] coordinates {(900,0)};
      \addplot[dashed, gray!70] coordinates {(300,0) (300,150)};
      \node[font=\scriptsize] at (axis cs:150,-45) {$a$};
      \node[font=\scriptsize] at (axis cs:370,80) {$b$};
      \node[font=\scriptsize] at (axis cs:610,-45) {$s=\rho-a$};
      \node[font=\scriptsize] at (axis cs:930,55) {$G$};
      \node[font=\scriptsize] at (axis cs:255,195) {$S_2$};
    \end{axis}
  \end{tikzpicture}
  \caption{一次检测后的可行集，与第二检测点的参数定义：$a$ 为沿示向线的前进量，
           $b$ 为侧移量，$s=\rho-a$ 为沿轴错位}
  \label{fig:p2set}
\end{figure}

\subsubsection{交会几何的闭式解与两次检测的极限}
\label{sec:p2geo}

沿用上节的局部标架与记号 $(a,b,s)$。将 $\sin\varphi=|b|/r_2$、$\cos\varphi=|s|/r_2$、
$r_2=\sqrt{s^2+b^2}$ 代入第 \ref{sec:p1} 节的直径近似式，得
\begin{equation}
  d(a,b;\rho)\approx\frac{2\tan\delta}{|b|}\sqrt{s^2+b^2}\,\sqrt{(\rho+|s|)^2+b^2}.
  \label{eq:p2d}
\end{equation}
式 \eqref{eq:p2d} 在 $b\to0$ 时发散，对应两检测点与干扰源共线、无法交会的退化情形。

\begin{theorem}[最优侧移与最小直径]
\label{thm:p2opt}
固定 $\rho$ 与前进量 $a$，式 \eqref{eq:p2d} 关于 $b$ 的极小点与极小值为
\begin{equation}
  b^{*}=\sqrt{|s|\,(\rho+|s|)},
  \qquad
  d_{\min}(a)=2\tan\delta\,\bigl(\rho+2|s|\bigr).
  \label{eq:p2bstar}
\end{equation}
\end{theorem}

\begin{proof}
对 $\ln d$ 求导并置零得 $-1/b+b/(s^2+b^2)+b/((\rho+|s|)^2+b^2)=0$，
整理为 $b^4=s^2(\rho+|s|)^2$，取正根即得 $b^{*}$。回代时注意
$s^2+b^{*2}=|s|(\rho+2|s|)$ 与 $(\rho+|s|)^2+b^{*2}=(\rho+|s|)(\rho+2|s|)$，
其乘积为 $b^{*2}(\rho+2|s|)^2$，开方后与 $2\tan\delta/b^{*}$ 相乘即得 $d_{\min}$。
\end{proof}

式 \eqref{eq:p2bstar} 的含义是清晰的：以正横站位（$s=0$，即 $\varphi=90^\circ$）的理想值
$2\rho\tan\delta$ 为基准，\textbf{沿轴每错位 $1$ 米，定位区域直径线性增加
$4\tan\delta\approx0.07$ 米}；例如 $|s|=0.5\rho$ 时直径已是理想值的 $2$ 倍。
换言之，距离估计的相对误差几乎按一比一转化为定位精度的损失。

式 \eqref{eq:p2bstar} 还给出一条与第二点位置\textbf{无关}的下界：
$d\ge d_{\min}(a)\ge 2\rho\tan\delta$。即无论第二检测点选在何处，
定位区域直径都不可能小于 $2\rho\tan\delta$，这是\textbf{两次检测所能达到的精度极限}。
结合一次清除的必要条件 $d\le2R_c$（因 $R_{\min}\ge d/2$）即得源距上限
\begin{equation}
  \rho\le R_c/\tan\delta=1145.8\ \text{米}.
  \label{eq:p2limit}
\end{equation}
设 $\rho$ 已知并穷举 $(a,b)$（附加“第二检测点与干扰源至少相距 $R_c$”这一现实约束，
否则应直接清除而非再次检测），所得最小 $R_{\min}$ 在 $\rho=1100,1146,1300$ 米处分别为
$19.21,20.01,21.97$ 米，与式 \eqref{eq:p2limit} 吻合。故有：

\begin{center}
\emph{源距超过 $1146$ 米时，即便事先已知距离，两次检测也不足以一次清除；
第二次检测应以将距离歧义由 $1495$ 米压缩至百米量级为目标。}
\end{center}

此结论与问题一定理 \ref{thm:clearable} 并不矛盾：后者设两检测点距源\textbf{均}不超过
$\ell$ 且张角 $90^\circ$，给出可主动构造的充分条件 $\ell^{*}\approx770$ 米；
而式 \eqref{eq:p2limit} 是第一检测点已由题目固定时的理论上限。
此外由第 \ref{sec:p1} 节，$\theta_1-\theta_2$ 恰等于交会角，故定位区域有界当且仅当
$\varphi>2\delta$，这是候选区域的\textbf{可行性门槛}，与“$\varphi$ 越接近 $90^\circ$ 越好”
这一优化方向属于不同层次。

\begin{theorem}[一致最优点的不存在性]
\label{thm:p2none}
不存在对全部 $\rho\in(5,\rho_{\mathrm{hi}}]$ 一致最优的第二检测点。
\end{theorem}

\begin{proof}
由式 \eqref{eq:p2bstar}，$d_{\min}(a)=2\tan\delta(\rho+2|\rho-a|)$ 关于 $a$ 的极小点集
恰为单点集 $\{a=\rho\}$。不同 $\rho$ 对应的极小点集互不相交，故不存在同时使所有 $\rho$
取到极小的 $a$。
\end{proof}

定理 \ref{thm:p2none} 表明“使定位区域直径最小的第二检测点”这一提法本身不适定，
必须代之以稳健化准则；而不同的稳健化方式会给出相距数百米的不同答案，
这正是下一节要解决的问题。

\subsubsection{代价加权与候选区域}
\label{sec:p2region}

\paragraph{行程即代价.}
题目在问题三中明确“派出\textbf{一只}机器狗”，附件 2 第 1.4 节亦规定每次测试只允许一个机器狗进入。
因此\textbf{两个检测点并非两台同时布设的传感器，而是同一条路径上的先后两站}：
第二检测点只能靠机器狗自行走到，而行进所消耗的正是题目要求最小化的时间。
脱离行程孤立地谈“较好的定位效果”，在数学上等价于假设存在第二台可瞬时布设的设备，与题设不符。

为量化这一代价，定义\textbf{绕行} $\Delta=|S_1S_2|+|S_2G|-\rho$，
即相对于沿示向线直趋干扰源多走的路程。按假设 (8) 的先验取期望（$\mathbb{E}[\rho]=855$ 米），
结果见表 \ref{tab:p2detour}：几何上最优的站位需多走近一公里，而本文推荐的站位几乎就在必经之路上。

\begin{table}[H]
  \centering
  \caption{不同第二检测点的期望绕行（$\mathbb{E}[\rho]=855$ 米）}
  \label{tab:p2detour}
  \begin{tabular}{lcccc}
    \toprule
    第二检测点 $(a,b)$ & $|S_1S_2|$/米 & $\mathbb{E}|S_2G|$/米 & 绕行 $\Delta$/米 & 折合/秒 \\
    \midrule
    本文推荐 $(450,120)$      & $466$  & $470$  & $\mathbf{80}$   & $16.0$ \\
    几何准则最优 $(900,580)$  & $1071$ & $657$  & $873$           & $174.5$ \\
    极小极大最优 $(1125,925)$ & $1456$ & $1009$ & $\mathbf{1610}$ & $322.1$ \\
    \bottomrule
  \end{tabular}
\end{table}

据此，本文以\textbf{期望总清除时间} $T$（自第一次检测起至该源被清除为止的虚拟时间，
含行进、检测与清除）为准则，目标为 $\min_{(a,b)}\mathbb{E}_\rho[T]$。
由机器狗速度 $5$ 米/秒与单次检测耗时 $5$ 秒，可将时间与路程统一到\textbf{等效行程}量纲：
\begin{equation}
  \text{一次检测}\;=\;5\ \text{秒}\;=\;25\ \text{米行程},
  \label{eq:p2exch}
\end{equation}
于是各项代价可直接相加比较。这一换算贯穿本节的全部权衡。

\paragraph{精度的边际收益递减.}
图 \ref{fig:p2front} 给出\textbf{效率前沿}：给定至第二检测点的行进距离，
期望定位区域直径 $\mathbb{E}[d]$ 所能达到的最小值。两个事实决定了本问的答案。

其一，$\mathbb{E}[d]$ 在 $50.2$ 米处饱和——在整个 $(a,b)$ 平面上取不到更小的值，
行进距离由 $1050$ 米再增加 $500$ 米反而回升至 $58.5$ 米。而一次清除要求 $d\le40$ 米，
\textbf{该下确界高于门限}。这意味着 $\mathbb{E}[d]=52$ 米与 $209$ 米在“是否仍需第三次检测”
上并无区别，为压低前者所增加的行程换取的是\textbf{超出门限、无法利用的精度}。

其二，$\mathbb{E}[d]$ 较优的点，一次清除概率未必高：$(1018,587)$ 处 $\mathbb{E}[d]=51.6$ 米
而一次清除概率为 $0$，$(700,430)$ 处 $\mathbb{E}[d]=57.8$ 米而该概率达 $0.445$。
其原因在于，行至千米之外后对任何源距第二楔形均较宽，直径处处约 $50$ 米、处处恰好不足；
而稍近的站位对部分源距可压入 $20$ 米以内。\textbf{对一个仅在门限处有意义的量取期望，
会被“处处恰好不足”的解所误导}，这是 $\mathbb{E}[d]$ 准则本身的局限。

\begin{figure}[H]
  \centering
  \begin{tikzpicture}
    \begin{axis}[
      width=0.62\linewidth, height=4.8cm, ymode=log,
      xlabel={至第二检测点的行进距离 / 米}, ylabel={$\mathbb{E}[d]$ / 米},
      xmin=80, xmax=1720, grid=major,
      ytick={40,50.2,70,100,150,200,300,500},
      yticklabels={40,50.2,70,100,150,200,300,500},
      minor ytick={}, log ticks with fixed point,
      tick label style={font=\scriptsize},
      legend style={font=\scriptsize, at={(0.97,0.97)}, anchor=north east, draw=black!30},
      tick label style={font=\scriptsize}, label style={font=\scriptsize},
    ]
      \addplot[blue!70!black, thick, mark=*, mark size=1.1pt]
        table [x=行进距离, y=几何最优_ED, col sep=comma] {\datadir/问题二_效率前沿.csv};
      \addlegendentry{给定行程下可达的 $\min\mathbb{E}[d]$}
      \addplot[black, dotted, thick, domain=80:1720] {50.2};
      \addlegendentry{下确界 $50.2$ 米}
      \addplot[red, thick, dashed, domain=80:1720] {40};
      \addlegendentry{一次清除门限 $2R_c=40$ 米}
    \end{axis}
  \end{tikzpicture}
  \caption{定位精度的效率前沿。边际收益迅速递减，且下确界 $50.2$ 米高于可清除门限 $40$ 米，
           故增加行程无法换来可用的精度}
  \label{fig:p2front}
\end{figure}

\paragraph{两个方向的权衡.}
按式 \eqref{eq:p2exch} 折算后，侧移量与前进量各自的权衡可直接相减比较。

侧移量 $b$ 的代价是绕行，当 $b\ll a$ 时约为 $b^2\rho/[2a(\rho-a)]$，是\textbf{二阶}量；
其收益是定位区域变小从而后续补测绕行减少，由式 \eqref{eq:p2d} 知正比于 $1/b$。
令代价合计为 $Ab^2+B/b$，极小点 $b^{*}=(B/2A)^{1/3}$，代入化简得 $b^{*}\approx0.10\rho$，
即\textbf{侧移量约取源距的十分之一}。注意此处为\textbf{立方根}依赖：
欲使最优侧移改变一倍需使代价比改变八倍，故目标函数在该方向上极不敏感——
这正是本问答案呈现为一个区域而非一个点的数学根源。

前进量 $a$ 的绕行 $\Delta(a)$ 关于 $a$ 呈 U 形：在起点侧移需先横向走出再斜向折回，
绕行为\textbf{一阶}量 $\approx b$；而在行程中段侧移时路径近乎直线，绕行降为二阶量
$\approx2b^2/\rho$（$b=120$、$\rho=855$ 时分别为 $128$ 米与 $34$ 米）。
另一侧的代价是\textbf{超越真实源距后需往返两遍}，而 $\rho$ 恰为未知量：
取 $a=300,450,600$ 米时超越概率分别为 $6\%,13\%,23\%$，期望往返惩罚分别为 $11,38,91$ 米。
由 $a$ 从 $300$ 增至 $600$ 米仅使后续绕行减少约 $23$ 米而惩罚增加 $80$ 米，
故最优前进量应取源距分布的一个\textbf{低分位点}而非中位数。

上述闭式分析给出的是量级判断：$b^{*}\approx0.10\rho$ 在 $\rho=855$ 米处约合 $86$ 米，
$a$ 的最优值落在 $300\sim450$ 米。它们与下文由全平面扫描得到的推荐点
$(450,120)$ 及候选区域 \eqref{eq:p2region} 一致，可视为对数值结果的解析印证；
闭式偏小是因为它以单一代表源距 $\mathbb{E}[\rho]$ 计算，而扫描是对整个先验取期望。

\paragraph{候选区域.}
在 $(a,b)$ 平面上以 $25\times15$ 米的步长扫描 $7120$ 个格点，
对每点在 $\rho$ 上取期望后分别计算两种准则，结果见图 \ref{fig:p2region}。
两准则的最优点相距约 $700$ 米；表 \ref{tab:p2compare} 进一步给出四种稳健化准则各自的最优点。

\begin{table}[H]
  \centering
  \caption{四种准则给出四个不同的最优第二检测点}
  \label{tab:p2compare}
  \begin{tabular}{llcccc}
    \toprule
    准则 & 最优点 $(a,b)$ & $\mathbb{E}[d]$/米 & $\max D$/米 & 一次清除概率 & $\mathbb{E}[T]$/秒 \\
    \midrule
    $\min\mathbb{E}[d]$（几何）    & $(900,580)$  & $\mathbf{50.2}$ & $92.8$ & $0.256$ & $363.5$ \\
    $\min\max D$（极小极大）        & $(1125,925)$ & $57.2$ & $\mathbf{75.9}$ & $0.000$ & $516.3$ \\
    $\max P(\text{一次清除})$       & $(700,430)$  & $57.8$ & --- & $\mathbf{0.445}$ & $291.0$ \\
    $\min\mathbb{E}[T]$（本文采用） & $(450,145)$  & $172.9$ & $398.6$ & $0.238$ & $\mathbf{224.1}$ \\
    \bottomrule
  \end{tabular}
\end{table}

表 \ref{tab:p2compare} 说明\textbf{“最优”离开准则没有意义}，准则的选取才是本问的核心。
本文依题目正文“此任务完成时间越短越好”与上述单机器狗论证，采用 $\min\mathbb{E}[T]$。
取相对最优值 $5\%$ 的容差，并在该水平集内取整体落入其中的内接矩形，得候选区域
\begin{equation}
  \boxed{\;
  a\in[0.10,\,0.30]\,\rho_{\mathrm{hi}},\qquad
  |b|\in[0.05,\,0.12]\,\rho_{\mathrm{hi}},\qquad\text{左右对称}
  \;}
  \label{eq:p2region}
\end{equation}
（$\rho_{\mathrm{hi}}=1500$ 米时即 $a\in[150,450]$ 米、$|b|\in[75,180]$ 米，
相当于沿与示向度成 $7^\circ\!\sim\!45^\circ$ 的方向行进 $230\!\sim\!690$ 米）。该区域满足三条性质：

\begin{enumerate}[label=(\roman*)]
  \item \textbf{整体位于接收可靠域内}。由 $\rho\le R_{\mathrm{eff}}$ 知
        $R_{\mathrm{eff}}\ge\max(1000,\rho)$，故 $|S_2G|\le\max(1000,\rho)$ 即可保证收到信号；
        对 $\rho$ 取遍可行区间，该条件的解为两圆盘之交 $D((5,0),1000)\cap D((1000,0),1000)$，
        即图 \ref{fig:p2region} 中的透镜形边界。$4\times10^5$ 次抽样表明域内接收率恒为 $1$。
  \item \textbf{满足有界性门槛}。对区域内全部点与 $\rho\in(5,1500]$ 全程扫描，
        最小交会角为 $3.18^\circ$，为门槛 $2\delta$ 的 $1.6$ 倍。
  \item \textbf{区域内各点的期望总清除时间均不超过最优值的 $1.05$ 倍}，
        见第 \ref{sec:p2check} 节的逐点检验。
\end{enumerate}

\begin{figure}[H]
  \centering
  \begin{tikzpicture}[
    /pgfplots/every axis/.append style={
      width=0.80\linewidth, height=4.6cm, view={0}{90},
      ylabel={侧移量 $b$ / 米},
      xmin=-200, xmax=1500, ymin=0, ymax=1000,
      colormap name=viridisrev, colorbar,
      colorbar style={font=\scriptsize, width=0.24cm},
      title style={font=\small},
      tick label style={font=\scriptsize}, label style={font=\scriptsize},
      xtick={0,250,500,750,1000,1250,1500}, ytick={0,250,500,750,1000},
    }]
    \begin{axis}[
      name=g, point meta min=50, point meta max=300,
      colorbar style={ylabel={$\mathbb{E}[d]$ / 米}},
      title={(a) 几何准则 $\min\mathbb{E}[d]$：最优点远在 $(900,580)$},
      xticklabels={},
    ]
      \addplot3[surf, shader=interp, mesh/rows=34]
        table [x=a, y=b, z=z, col sep=comma] {\datadir/问题二_等值面直径.csv};
      \addplot3[white, very thick] coordinates
        {(750,400,1e4) (1050,400,1e4) (1050,835,1e4) (750,835,1e4) (750,400,1e4)};
      \addplot3[only marks, mark=star, mark size=4pt, draw=black, fill=white]
        coordinates {(900,580,1e4)};
      \node[font=\scriptsize, black] at (axis cs:900,320,1e4) {几何准则的 $5\%$ 域};
    \end{axis}
    \begin{axis}[
      name=t, at={($(g.south)-(0,1.1cm)$)}, anchor=north,
      point meta min=216, point meta max=420,
      colorbar style={ylabel={$\mathbb{E}[T]$ / 秒}},
      title={(b) 时间准则 $\min\mathbb{E}[T]$（本文采用）：最优点移至 $(450,145)$},
      xlabel={前进量 $a$ / 米},
    ]
      \addplot3[surf, shader=interp, mesh/rows=34]
        table [x=a, y=b, z=z, col sep=comma] {\datadir/问题二_等值面时间.csv};
      \addplot3[green!40!black, thick, dashed] table [x=x, y=y, z expr=1e4, col sep=comma]
        {\datadir/问题二_接收可靠域.csv};
      \addplot3[red, very thick] coordinates
        {(150,75,1e4) (450,75,1e4) (450,180,1e4) (150,180,1e4) (150,75,1e4)};
      \addplot3[only marks, mark=star, mark size=4pt, draw=black, fill=red]
        coordinates {(450,120,1e4)};
      \addplot3[only marks, mark=o, mark size=3pt, thick, white]
        coordinates {(900,580,1e4)};
      \node[font=\scriptsize, black] at (axis cs:300,300,1e4) {候选区域};
      \node[font=\scriptsize, green!55!white] at (axis cs:330,760,1e4) {接收可靠域边界};
      \node[font=\scriptsize, white] at (axis cs:900,700,1e4) {几何准则最优点};
    \end{axis}
  \end{tikzpicture}
  \caption{两种准则下的目标函数等值面（色标已反转，亮色为优）。
           几何准则的最优点位于 $(900,580)$，时间准则的最优点移至 $(450,145)$，
           二者相距约 $700$ 米；红色矩形为式 \eqref{eq:p2region} 给出的候选区域}
  \label{fig:p2region}
\end{figure}

\paragraph{检测策略.}
式 \eqref{eq:p2region} 按 $\rho_{\mathrm{hi}}$ 缩放，故对任意起测点与示向度自动适用。
完整策略见算法 \ref{alg:p2}，其中后续补测沿用同一原则：新楔形应横切当前定位区域的长轴，
侧移量取该区域半长的 $\kappa$ 倍。由式 \eqref{eq:p2exch}，“单步到位”（$\kappa=1$）
虽可一次将区域压缩至极小却需多走一个半长；实测 $\kappa=1$ 时平均耗时 $227.3$ 秒，
$\kappa=0.25$ 时为 $217.0$ 秒，差额 $10$ 秒恰合 $50$ 米、约两次检测的代价，故取 $\kappa=0.25$。

\begin{algorithm}[H]
  \caption{基于集员估计与代价加权的单源定位清除策略}
  \label{alg:p2}
  \begin{algorithmic}[1]
    \Require 第一检测点 $S_1$ 与示向度 $\theta_1$；参数 $\lambda=0.30$、$\mu=0.08$、$\kappa=0.25$
    \Ensure 干扰源被清除
    \State $\rho_{\mathrm{hi}}\gets\min(1500,\ \text{沿楔形至目标圆边界的距离})$
    \State $\Omega\gets W(S_1,\theta_1,\delta)\cap\{5<\rho\le\rho_{\mathrm{hi}}\}\cap D(O,R)$
    \State 取 $\sigma\in\{+1,-1\}$ \Comment{侧移方向，问题二中任选，问题三中取朝向其余待定位源的一侧}
    \State $P\gets S_1+\lambda\rho_{\mathrm{hi}}\,u(\theta_1)+\sigma\mu\rho_{\mathrm{hi}}\,n(\theta_1)$
    \While{真}
      \State 移动至 $P$ 并检测，得结果 $z$
      \If{$z=\texttt{near}$} \State 就地清除；\textbf{返回} \EndIf
      \If{$z=\texttt{direction}$} \State $\Omega\gets\Omega\cap W(P,\theta,\delta)$ \EndIf
      \State 计算 $\Omega$ 的最小覆盖圆 $(c,R_{\min})$ 与直径 $d$
      \If{$R_{\min}\le R_c$} \State 移动至 $c$ 并清除；\textbf{返回} \EndIf
      \State $n\gets\Omega$ 直径方向的法向（取靠近当前位置的一侧）
      \State $P\gets c+\kappa\,(d/2)\,n$ \Comment{在区域中心的正横方向补测，横切长轴}
    \EndWhile
  \end{algorithmic}
\end{algorithm}

\subsubsection{数值结果与模型检验}
\label{sec:p2check}

\paragraph{仿真设置.}
按附件 1 表 1 与附件 2 第 4 节复刻模拟器的计时规则（移动按距离除以速度、检测 $5$ 秒、
切换频道 $1$ 秒、清除成功 $5$ 秒、未发现 $3$ 秒），并将示向度误差实现为\textbf{位置的确定性函数}
以复刻“同一地点重复检测不改变误差”，示向度截取两位小数、算法侧楔形半角取 $1.005^\circ$。
按假设 (8) 抽样并仅保留在 $S_1$ 能收到信号的算例，共 $300$ 组，随机种子 $20260910$。

\paragraph{与理论下界的比较.}
设机器狗一开始即已知干扰源位置，直趋并清除，耗时为“源距除以速度 $5$ 米/秒，
再加清除的 $5$ 秒”，此为\textbf{任何策略的下界}，
与所用算法无关。本批算例平均源距 $844.7$ 米，故下界为 $173.9$ 秒。各策略对比见表 \ref{tab:p2compare2}。

\begin{table}[H]
  \centering
  \caption{各策略的平均总清除时间（同一批 $300$ 组算例）}
  \label{tab:p2compare2}
  \begin{tabular}{lccc}
    \toprule
    策略 & 平均总清除时间/秒 & 相对下界 & 平均检测次数 \\
    \midrule
    理论下界（已知干扰源位置）              & $173.9$ & $1.00$ & --- \\
    \textbf{本文策略}                       & $\mathbf{216.5}$ & $\mathbf{1.24}$ & $2.96$ \\
    微小位移（$\lambda=\mu=0.02$）          & $246.2$ & $1.42$ & $3.49$ \\
    单次交会后覆盖清除                      & $251.7$ & $1.45$ & $2.00$ \\
    纯径向推进（示向度反转二分）            & $257.7$ & $1.48$ & $3.10$ \\
    覆盖式清除搜索（不作交会）              & $302.4$ & $1.74$ & $1.00$ \\
    几何准则最优点 $(900,580)$              & $366.1$ & $2.10$ & $2.75$ \\
    纯横向偏移（$\lambda=0,\mu=0.5$）       & $409.2$ & $2.35$ & $3.53$ \\
    极小极大准则最优点 $(1125,925)$         & $516.3$ & $2.97$ & $3.09$ \\
    \bottomrule
  \end{tabular}
\end{table}

本文策略的平均总清除时间为理论下界的 $1.24$ 倍，故\textbf{任何策略相对本文策略的改进幅度
不超过 $19.7\%$}，且该上界只有在事先已知干扰源位置时方能达到；次优方案已落后 $14\%$。
表中“覆盖式清除搜索”指不作交会定位而以 $81$ 个清除点覆盖整条楔形
（构造经 $2\times10^5$ 次抽样验证零遗漏）：因清除“未发现”仅耗时 $3$ 秒而短于一次检测，
该方案构成实质竞争，但实测仍慢 $40\%$。
按式 \eqref{eq:p2exch} 折算可见代价去向：移动占 $91\%$、检测占 $7\%$、清除占 $2\%$；
超出理论下界的部分中，绕行占 $65\%$、额外检测占 $23\%$，两者量级相当，是均衡最优的特征。

\paragraph{候选区域的逐点检验.}
式 \eqref{eq:p2region} 并非事后描述，而是以“区域内各点的实测平均耗时不超过最优值 $1.05$ 倍”
为\textbf{定义}。对其四角、中心与推荐点逐点仿真，结果见表 \ref{tab:p2verify}：
区域内全部检验点均不超过 $1.046$，区域外各点显著劣化。
需说明的是，$5\%$ 水平集本身大于该矩形且非矩形（例如 $a=100$、$a=600$ 与 $b=35$ 处实测均约 $1.036$），
采用内接矩形是为便于表述与使用，其代价是结论偏保守：\textbf{区域外未必不佳，而区域内必定达标}。

\begin{table}[H]
  \centering
  \caption{候选区域的逐点检验（$\rho_{\mathrm{hi}}=1500$ 米）}
  \label{tab:p2verify}
  \begin{tabular}{lccccc}
    \toprule
    & \multicolumn{2}{c}{区域四角} & 区域中心 & 推荐点 & 区域外 \\
    \cmidrule(lr){2-3}\cmidrule(lr){4-4}\cmidrule(lr){5-5}\cmidrule(lr){6-6}
    $(a,b)$/米 & $(150,75)$ & $(150,180)$ & $(300,127.5)$ & $(450,120)$ & $(800,120)$ \\
    平均耗时/秒 & $222.1$ & $226.5$ & $216.0$ & $216.5$ & $249.0$ \\
    相对最优    & $1.026$ & $\mathbf{1.046}$ & $0.998$ & $1.000$ & $1.150$ \\
    \midrule
    $(a,b)$/米 & $(450,75)$ & $(450,180)$ & --- & --- & $(900,580)$ \\
    平均耗时/秒 & $216.5$ & $222.1$ & --- & --- & $366.1$ \\
    相对最优    & $1.000$ & $1.026$ & --- & --- & $1.691$ \\
    \bottomrule
  \end{tabular}
\end{table}

\paragraph{模型检验.}
本文的最优点由两条互不共享推导的途径得到：基于闭式收尾代价的代理模型全平面扫描给出 $(450,145)$，
端到端仿真的参数搜索给出 $(450,120)$；将前者代入仿真实测为 $217.9$ 秒，仅劣 $0.6\%$，
二者落在同一平坦区内。

几何实现按第 \ref{sec:p1check} 节的原则复核，复核代码不调用被检验代码，而由题面定义重新推导：
楔形隶属按“方位角与示向度之差不超过 $\delta$”直接判定；定位区域顶点由边界曲线两两求交后
按原始定义筛选；凸包、直径与最小覆盖圆均独立实现。结果为：$600$ 组随机构型的直径最大偏差
$9.0\times10^{-3}$ 米（源于外圈圆弧的折线化，且模型侧恒不小于复核侧，偏于安全）；
$10^6$ 个随机点的隶属判定中“依定义属于而模型判定不属于”的个数为 $0$；
最小覆盖圆除与问题一的枚举算法比对外，另以\textbf{最优性证书}直接验证——
即检验“全部顶点被覆盖”与“圆心落在支撑点集的凸包内”这一对充要条件\cite{rockafellar}，
后者保证半径不可再缩，$3000$ 个随机凸多边形全部通过。
变异测试向模型注入楔形半角错误、顶点偏移、
遗漏距离约束等六类缺陷，复核\textbf{全部检出}。

\paragraph{对假设的灵敏度.}
本节的全部解析结论（式 \eqref{eq:p2d}--\eqref{eq:p2limit}、引理 \ref{lem:p2para}、
定理 \ref{thm:p2opt} 与 \ref{thm:p2none}）只依赖确界 $|\varepsilon|\le\delta$，
与误差的分布形态及相关结构无关。数值结果同样稳健：将误差改为三角分布、取最坏值 $\pm\delta$、
或引入相关长度 $300$ 至 $1000$ 米的空间相关误差场，平均耗时的变动分别为
$-0.6\%$、$+2.0\%$、$-0.7\%\!\sim\!-0.2\%$，而“定位区域包含干扰源真实位置”的比例
在所有情景下恒为 $100\%$——这是集员估计的固有保证。

至于先验假设 (8)，将其换为 $\rho$ 上的均匀分布后最优侧移量几乎不变（$145\to130$ 米），
而最优前进量由 $450$ 米移至 $225$ 米。故“侧移量取源距的十分之一”这一结论稳健，
而前进量对先验较为敏感，使用时应以式 \eqref{eq:p2region} 的区间而非单点为准。
此外，以 $(0,0)$、$(900,0)$、$(1400,0)$、$(1000,1000)$ 为起测点分别重新抽样并就地重搜参数，
沿用固定参数 $(\lambda,\mu)=(0.30,0.08)$ 相对于就地重搜的劣化均不超过 $0.4\%$，
表明按 $\rho_{\mathrm{hi}}$ 缩放的自适应性成立。

\subsection{问题三：全向干扰源的自动搜索定位与清除策略}

% TODO: 待完成。含 5.3.1 模型建立、5.3.2 算法设计、5.3.3 演练与正式测试结果、5.3.4 模型检验。

\subsection{问题四：含定向干扰源时的检测与清除策略}

% TODO: 待完成。含 5.4.1 模型建立、5.4.2 算法设计、5.4.3 测试结果、5.4.4 模型检验。

\section{模型的评价与推广}

\subsection{模型的优点}

\begin{enumerate}[label=(\arabic*)]
  \item \textbf{与题目的误差信息形式相匹配。} 题目给出的是误差确界而非分布，
        且明确指出同一地点重复检测不改变误差。本文采用的集员估计方法只依赖确界，
        所得定位区域是严格包含真实位置的集合，不需要任何分布假设。
  \item \textbf{算法结构简单且退化情形完备。} 楔形化为半平面之交后，
        任意 $n$ 的处理方式统一；表 \ref{tab:degenerate} 给出了五种形态的完整分类，
        并对空集与无界情形给出明确的检出机制，不会输出没有实际意义的结果。
  \item \textbf{结论经独立方法交叉验证。} 第 \ref{sec:p1check}、\ref{sec:p2check} 节的检验
        均不调用被检验代码，且经变异测试确认其检错能力。
  \item \textbf{评价准则由题意而非习惯确定。} 问题二中，本文没有沿用“定位区域越小越好”的
        直观准则，而是先证明该提法不适定（定理 \ref{thm:p2none}），再依据“只派出一只机器狗”
        与“完成时间越短越好”两条题设，改用期望总清除时间为准则；
        表 \ref{tab:p2compare} 表明不同准则给出的最优点相距数百米，准则的选取确有实质影响。
  \item \textbf{给出了最优性的可验证上界。} 以“已知干扰源位置”构造理论下界，
        得出本文策略的可改进幅度不超过 $19.7\%$（表 \ref{tab:p2compare2}）。
        该结论不依赖于与其他具体方法的比较，因而不会因遗漏某种竞争方案而失效。
\end{enumerate}

\subsection{模型的缺点}

\begin{enumerate}[label=(\arabic*)]
  \item 集员估计给出的是包含真实位置的区域，\textbf{不提供区域内的概率分布}。
        若后续需要按概率加权决策，需另行引入分布假设。
  \item 问题一的定位区域\textbf{未并入}“干扰源必位于目标区域内”这一先验：
        引入后当干扰源接近边界时会产生圆弧边，使定位区域不再是多边形，
        与题目“多边形定位区域”的措辞不符。经测算此情形仅出现于最外约 $2\%$ 的面积范围内，
        对直径的影响在厘米量级，故所得直径是偏保守的上界。
        需要说明的是，问题二的可行集 $\Omega_1$（式 \eqref{eq:p2omega}）\textbf{已并入}该先验
        以及距离约束，因为该问不要求输出多边形，且这两条约束对候选区域的位置有实质影响。
  \item 问题二的策略在一个\textbf{参数族}内取优，而非全部策略中的最优。
        真正的最优需求解部分可观测的序贯决策问题，本文未予求解；
        其影响已由理论下界限定在 $19.7\%$ 以内（表 \ref{tab:p2compare2}）。
  \item 问题二的期望值依赖于干扰源位置的先验假设 (8)。灵敏度分析表明侧移量的结论稳健，
        而前进量对先验较为敏感，故本文以区间而非单点给出候选区域。
\end{enumerate}

\subsection{模型的推广}

问题一建立的定位区域模型与直径算法是后续三问的共同基础：问题二在其上叠加距离约束
构成可行集，并以 $R_{\min}\le R_c$ 作为可实施清除的判据；问题三、四中，
每次新增检测只需对当前可行集再作一次半平面裁剪，即可增量更新，无需重算。

问题二确立的\textbf{代价加权}思路对问题三、四尤为关键：由于只有一只机器狗，
检测点的选取与路径规划实为同一问题，“定位精度”与“行进代价”必须在同一量纲下权衡。
式 \eqref{eq:p2exch} 给出的换算（一次检测等价于 $25$ 米行程）可直接沿用；
而候选区域式 \eqref{eq:p2region} 是一个区域而非一个点，这一自由度使得在问题三中
可以选取同时落入多个干扰源候选区域的位置，实现\textbf{一次停靠服务多个频道}，
从而把检测成本摊薄。

本文方法可推广至一般的\textbf{仅测向、不测距}的定位问题：只要观测误差以确界形式给出、
且观测代价与观测位置相关，本文“集员估计 $+$ 代价加权”的框架即可适用。

% ==========================================================
% AI 工具使用声明
%
% 官方《人工智能工具使用规定（2026 年试行）》要求本声明置于参考文献之前，
% 下面两条命令二选一，不可省略:
%
%   \CumcmAINotUsed          未使用任何 AI 工具
%   \CumcmAIUsed{简要用途}    使用了 AI 工具, 需在支撑材料中附
%                            「AI工具使用详情.pdf」
% ==========================================================
\CumcmAIUsed{代码编写与调试、数值实验脚本、文档排版}

\begin{thebibliography}{9}
  \bibitem{schweppe} Schweppe F C. Recursive state estimation: Unknown but bounded errors
    and system inputs[J]. IEEE Transactions on Automatic Control, 1968, 13(1): 22--28.
  \bibitem{jung} Jung H W E. Über die kleinste Kugel, die eine räumliche Figur einschließt[J].
    Journal für die reine und angewandte Mathematik, 1901, 123: 241--257.
  \bibitem{sh} Sutherland I E, Hodgman G W. Reentrant polygon clipping[J].
    Communications of the ACM, 1974, 17(1): 32--42.
  \bibitem{rockafellar} Rockafellar R T. Convex Analysis[M]. Princeton: Princeton University Press, 1970.
  \bibitem{qhull} Barber C B, Dobkin D P, Huhdanpaa H. The quickhull algorithm for convex hulls[J].
    ACM Transactions on Mathematical Software, 1996, 22(4): 469--483.
\end{thebibliography}

\begin{appendices}

\section{问题一的定理证明}
\label{app:proof}

\paragraph{有界性判据.}
$L\ne\varnothing$ 时，$L$ 有界当且仅当 $\bigcap_{i=1}^{n}[\theta_i-\delta,\theta_i+\delta]=\varnothing$
（视作单位圆周上的弧）。事实上，对非空闭凸集族，交集的回收锥等于各回收锥之交，
即 $\operatorname{recc}(L)=\bigcap_i\operatorname{recc}(W_i)$\cite{rockafellar}；
楔形 $W_i$ 的回收锥恰为方向区间 $[\theta_i-\delta,\theta_i+\delta]$ 张成的锥。
非空闭凸集有界当且仅当其回收锥仅含零向量，而诸锥之交仅含零向量当且仅当对应方向区间之交为空。
该判据为 $O(n)$，且自动涵盖 $n=1$：单个区间与自身相交非空，直接判为无界。

\paragraph{定理 \ref{thm:counter} 步骤 (4)：直径由对径顶点对取到.}
六顶点方向相隔 $60^\circ$ 且半径交替，故对径顶点半径必不相同，其距离为 $r_1+r_2$。
其余顶点对中，相隔 $120^\circ$ 的同半径对距离为 $\sqrt3\,r$（$r\le r_2$），
相隔 $60^\circ$ 的异半径对距离为 $\sqrt{r_1^2+r_2^2-r_1r_2}<r_1+r_2$。
而 $r_1+r_2>\sqrt3\,r_2$ 当且仅当 $r_1/r_2>\sqrt3-1$，即
$\cos(30^\circ+\delta)/\cos(30^\circ-\delta)>\sqrt3-1$；在 $\delta=15^\circ$ 处取等号
（此时 $\cos45^\circ/\cos15^\circ=\sqrt3-1$ 精确成立）。
故对 $\delta\in(0^\circ,15^\circ]$ 有 $d=r_1+r_2$。

\paragraph{定理 \ref{thm:counter}(c)：覆盖比的闭式.}
三个半径为 $r_2$ 的顶点互成 $120^\circ$，构成以 $G$ 为外心、$r_2$ 为外接半径的正三角形，
其最小覆盖圆即该外接圆；半径 $r_1<r_2$ 的三个顶点位于其内部，故 $R_{\min}=r_2$。
利用和差化积 $\cos(30^\circ+\delta)+\cos(30^\circ-\delta)=\sqrt3\cos\delta$，得

\[
  \frac{R_{\min}}{d/2}=\frac{2r_2}{r_1+r_2}
  =\frac{2\cos(30^\circ-\delta)}{\sqrt3\cos\delta}
  =\frac{\sqrt3\cos\delta+\sin\delta}{\sqrt3\cos\delta}
  =1+\frac{\tan\delta}{\sqrt3},\qquad \delta\in(0^\circ,15^\circ].
\]

代入 $\delta=15^\circ$ 并注意 $\tan15^\circ=2-\sqrt3$，即得比值 $2/\sqrt3$，
恰为 Jung 界。故该界在“楔形之交”这一类区域内可以精确取到。
本题 $\delta=1^\circ$ 时该对称构型的比值为 $1.0101$，为 Jung 界的 $87.5\%$。

\section{核心代码}

\begin{lstlisting}[language=Python, caption={楔形化为半平面并求交（问题一）}]
def 楔形半平面(S, theta_deg, delta_deg):
    “”“顶点 S、中心方向 theta、半角 delta 的楔形，返回其两个半平面 (a, b, c)，
    含义为 a*x + b*y <= c。”“”
    lo = np.radians(theta_deg - delta_deg)
    hi = np.radians(theta_deg + delta_deg)
    u_lo = (np.cos(lo), np.sin(lo))
    u_hi = (np.cos(hi), np.sin(hi))
    h1 = (u_lo[1], -u_lo[0], u_lo[1] * S[0] - u_lo[0] * S[1])
    h2 = (-u_hi[1], u_hi[0], -u_hi[1] * S[0] + u_hi[0] * S[1])
    return [h1, h2]


def 求定位区域(检测点, 示向度, delta_deg):
    “”“n 个楔形取交，返回 (顶点表, 状态)。”“”
    if not 角区间交为空(示向度, delta_deg):
        return [], “无界”
    多边形 = [(-BOX, -BOX), (BOX, -BOX), (BOX, BOX), (-BOX, BOX)]
    for S, th in zip(检测点, 示向度):
        for a, b, c in 楔形半平面(S, th, delta_deg):
            多边形 = 裁剪(多边形, a, b, c)
            if not 多边形:
                return [], “空”
    return 去冗余(多边形), “多边形”
\end{lstlisting}

\begin{lstlisting}[language=Python, caption={直径与覆盖判定（问题一）}]
def 求直径(顶点):
    “”“凸多边形的直径必在顶点对取到（定理 2）。”“”
    最大, 点对 = 0.0, None
    for p, q in combinations(顶点, 2):
        d = float(np.hypot(p[0] - q[0], p[1] - q[1]))
        if d > 最大:
            最大, 点对 = d, (p, q)
    return 最大, 点对


def 覆盖判定(顶点, A, B):
    “”“泰勒斯判据：V 在圆内或圆上 <=> (A-V)·(B-V) <= 0。”“”
    A, B = np.asarray(A), np.asarray(B)
    for V in 顶点:
        if float(np.dot(A - np.asarray(V), B - np.asarray(V))) > 1e-7:
            return False
    return True
\end{lstlisting}


\begin{lstlisting}[language=Python, caption={一次检测后的可行集与最优侧移（问题二）}]
def 先验区域(S1, theta1, delta, rho_lo=5.0, rho_hi=None, R=1800.0, 弧段数=6):
    """可行集 Omega_1 = 楔形 ∩ {rho_lo < rho <= rho_hi} ∩ 目标圆盘。
    外圈圆弧用切线外逼近，保证所得区域是真实可行集的超集（宁可高估不可漏）。"""
    if rho_hi is None:
        rho_hi = 距离上界(S1, theta1, delta, R)
    多边形 = [(-BOX, -BOX), (BOX, -BOX), (BOX, BOX), (-BOX, BOX)]
    半平面 = list(楔形半平面(S1, theta1, delta))
    半平面 += 圆的外切半平面(S1, rho_hi, 弧段数, (theta1 - delta, theta1 + delta))
    u = np.array([np.cos(np.radians(theta1)), np.sin(np.radians(theta1))])
    半平面.append((-u[0], -u[1], -(rho_lo * np.cos(np.radians(delta)) + float(S1 @ u))))
    半平面 += 圆的外切半平面((0.0, 0.0), R, 360)
    for a, b, c in 半平面:
        多边形 = 裁剪(多边形, a, b, c)
        if not 多边形:
            return []
    return 去冗余(多边形)


def 最优侧移(rho, a):
    """定理 3：b* = sqrt(|s|(rho+|s|))，s = rho - a。"""
    s = np.abs(rho - a)
    return np.sqrt(s * (rho + s))


def 最优直径(rho, a, delta):
    """定理 3：d_min = 2 tan(delta) (rho + 2|s|)。"""
    s = np.abs(rho - a)
    return 2.0 * np.tan(np.radians(delta)) * (rho + 2.0 * s)
\end{lstlisting}

\begin{lstlisting}[language=Python, caption={单源定位清除策略（问题二，对应算法 2）}]
def 跑一局(sim, S1, theta1, LAM=0.30, MU=0.08, KAP=0.25, 频道=1, 最多检测=12):
    """集员估计 + 代价加权的单源策略。sim 为复刻模拟器计时规则的本地仿真器。"""
    rho_hi = 距离上界(S1, theta1, DELTA_算法)
    Omega = 先验区域(S1, theta1, DELTA_算法, rho_hi=rho_hi)
    目标 = 局部转全局((LAM * rho_hi, MU * rho_hi), S1, theta1)   # 侧移方向可取反
    for _ in range(最多检测):
        回 = sim.measure(目标, 频道)
        if 回["measure_result"] == "near":
            sim.clear(sim.位置, 频道)
            return
        if 回["measure_result"] == "direction":
            新 = 定位区域(Omega, 目标, 回["svd_deg"], DELTA_算法)
            if 新:
                Omega = 新
        d, R = 区域指标(Omega)[:2]
        c = np.asarray(最小覆盖圆快(Omega)[0], dtype=float)
        if R <= R_CLEAR:                      # 可一次清除，走到圆心清除
            sim.clear(c, 频道)
            return
        目标 = 补测点(Omega, sim.位置, KAP)    # 在区域正横方向补测，横切长轴


def 补测点(Omega, 位置, kappa):
    """新楔形应横切当前区域长轴；侧移量取区域半长的 kappa 倍。
    kappa 由式 (15) 的等效行程换算权衡得出，取 0.25。"""
    D, (A, B) = 求直径快(Omega)
    c = np.asarray(最小覆盖圆快(Omega)[0], dtype=float)
    轴 = np.asarray(B, float) - np.asarray(A, float)
    轴 = 轴 / max(float(np.hypot(*轴)), 1e-12)
    n = np.array([-轴[1], 轴[0]])
    if float(n @ (np.asarray(位置, float) - c)) < 0:
        n = -n
    return c + max(kappa * D / 2.0, 10.0) * n
\end{lstlisting}

\end{appendices}

\end{document}
